\begin{center}
	\Huge
	Bestemmelse af afledede funktioner
\end{center}

\section*{Den afledede for $x^2$}
\stepcounter{section}

Vi har set, hvordan differentialkvotienten for $x^2$ i punktet $x_0 = 2$ var lig $4$ ved at anvende definitionen for differentialkvotienten på funktionen i punktet. Vi vil nu vise, at den afledede funktion 
for $x^2$ er $2x$.

\begin{setn}
	Lad $f$ være givet ved
	\begin{align*}
		f(x) = x^2.
	\end{align*}
	Så gælder der, at den afledede funktion $f'$ er givet ved
	\begin{align*}
		f'(x) = 2x.
	\end{align*}
\end{setn}
\begin{proof}
	Vi vælger vores $x_0$ vilkårligt og opskriver sekanthældningen
	\begin{align*}
		a_s &= \frac{f(x_0+h)-f(x_0)}{h}\\
		&= \frac{(x_0+h)^2-x_0^2}{h}\\
		&\stackrel{(1)}{=} \frac{x_0^2+2x_0h+h^2-x_0^2}{h}\\
		&=\frac{2x_0h+h^2}{h}\\
		&=\frac{h(2x_0+h)}{h}\\
		&=2x_0+h,
	\end{align*}
	hvor vi i $(1)$ har anvendt kvadratsætningen $(a+b)^2 = a^2+b^2+2ab$. 
	
	Vi kan nu lade $h \to 0$. 
	\begin{align*}
		\lim_{h\to 0}a_s = \lim_{h\to 0}2x_0+h = 2x_0
	\end{align*}
	Da vi har valgt vores $x_0$ vilkårligt får vi derfor den afledede funktion
	\begin{align*}
		f'(x) = 2x.
	\end{align*}
\end{proof}

Vi husker på tabellen fra sidst
\begin{table}[H]
	\centering
	\rowcolors{1}{white}{teal!30}
	\begin{tabular}{cc}
		$f(x)$ & $f'(x)$ \\
		$k$ & $0$ \\
		$x$ & $1$ \\
		$ax + b$  & $a$ \\
		$x^2$ & $2x$ \\
		$x^a$ & $ax^{a-1} $
	\end{tabular}
	\caption{Udvalgte afledede funktioner}
	\label{tab:afledt}
\end{table}

\subsection*{Opgave 1}
\begin{enumerate}[label=\roman*)]
	\item Anvend definitionen af differentialkvotienten til at bestemme den afledede funktion for $f(x) = 2x^2$.
	\item Anvend definitionen af differentialkvotienten til at bestemme den afledede funktion for $f(x) = 4x$.
	\item Anvend definitionen af differentialkvotienten til at bestemme den afledede funktion for $f(x) = 3x^2 - 5x$.
	\item Anvend definitionen af differentialkvotienten til at bestemme den afledede funktion for $f(x) = 5x - 2$.
	\item Anvend definitionen af differentialkvotienten til at bestemme den afledede funktion for $f(x) = -4x^2 + 3x - 1$.
\end{enumerate}

\subsection*{Opgave 2}
\begin{enumerate}[label=\roman*)]
	\item Vis, at der for konstanter $k$ gælder, at $(k)' = 0$ ved at anvende definitionen af differentialkvotienten.	
	\item Vis, at $f(x) = ax + b$ har differentialkvotienten $f'(x) = a$ ved at anvende definitionen af differentialkvotienten.
\end{enumerate}


\subsection*{Opgave 3}
Bestem den afledede funktion for følgende polynomier ved at bruge regnereglerne for differentiation.
\begin{align*}
	&1) \ p(x) = \frac{1}{2}x^2 -\frac{5}{4}x^4   &&2) \ p(x) = 4x^{10} - 0.6x^{3}+\frac{2}{3}x^{2}+2x+1\\
	&3) \ p(x) = x^{1} + x^0  &&4) \ p(x) = 2x + 7     \\ 
\end{align*}

\subsection*{Opgave 4}
Differentiér følgende funktioner ved at bruge regnereglerne for differentiation.
\begin{align*}
	&1) \ f(x) = -\frac{4}{7}x^{\frac{2}{5}}   &&2) \ f(x) = 2x^{1.5}-0.5x^{-3}\\
	&3) \ f(x) = \frac{1}{x^2}  &&4) \ f(x) = 2\sqrt[3]{x}     \\ 
\end{align*}

\subsection*{Opgave 5}
\begin{enumerate}[label=\roman*)]
	\item Bestem differentialkvotienten for $f(x) = x^5+2x -3$ i punktet $x_0 = -2$
	\item Bestem differentialkvotienten for $g(x) = \frac{4}{5}x^{1.5}$ i punktet $x_0 =16$
	\item Bestem tangenthældningen for $h(x) = 2x^{2}-5x+10$ i punktet $x_0 = -4$
\end{enumerate}

\subsection*{Opgave 6}
Vis følgende differentialkvotienter ved at anvende definitionen af differentialkvotienten.
\begin{enumerate}[label=\roman*)]
	\item $f(x) = ax^2 + bx + c$ har $f'(x) = 2ax + b$.
	\item $f(x) = x^3$ har $f'(x) = 3x^2$.
\end{enumerate}